\section{Discussion}

\noindent\textbf{Architectural Efficiency and Benchmark Dynamics}: The experimental results demonstrate that sub-1M parameter architectures provide an effective operating point for clinical neuro-oncology, balancing high diagnostic fidelity with low edge latency. On standard multi-class classification, compact convolutional models with parallel atrous topologies achieve diagnostic parity with heavyweight backbones an order of magnitude larger, demonstrating that excessive parameter capacity yields diminishing returns on canonical cohorts. When evaluated on multi-sequence clinical glioma grading under cross-institutional acquisition shifts, compact models continue to show strong representation capacity, with Base LS-Net outperforming mobile baselines including ResNet-18 and MobileNet-V2. These findings confirm that tailored multi-scale inductive biases can capture complex neuropathological patterns under strict parameter constraints.

\noindent\textbf{Adaptive Test-Time Compute Allocation}: MCG-TTA establishes a principled framework for dynamic inference-time compute allocation. Rather than treating all incoming MRI slices uniformly with compute-heavy multi-view ensembles, MCG-TTA couples predictive confidence with multi-scale consistency verification. By deploying dual scale probes aligned with the internal dilation offsets of the atrous backbone, the method selectively interrogates borderline predictions where spatial sampling misalignment is most acute. This selective gating preserves single-pass baseline latency for confident inputs while effectively stabilizing ambiguous lesion boundaries. Crucially, by leaving high-confidence samples unperturbed, MCG-TTA prevents the diagnostic degradation associated with brute-force augmentation, maintaining high sensitivity for critical low-grade glioma cases.

\noindent\textbf{Methodological and Clinical Limitations}: Several limitations of the current study should be noted. First, evaluation was conducted on retrospective public cohorts (Nickparvar and BraTS 2020); prospective multi-center clinical validation is essential to verify performance across diverse scanner models and real-time hospital workflows. Second, the current inference pipeline evaluates diagnostic axial slices containing core lesion volumes; integrating an automated whole-brain slice filtering or segmentation stage would enable end-to-end volumetric screening from raw DICOM series. Third, hardware execution latency was profiled on standardized edge-representative CPU processors; evaluating on diverse embedded platforms, such as specialized neural processing units (NPUs) and mobile intra-operative consoles, represents an important next step.

\noindent\textbf{Future Directions for Edge Neuro-Oncology}: These directions suggest clear pathways for future research: conducting prospective multi-site clinical trials, integrating compact classification with real-time tumor boundary segmentation, exploring hardware-specific kernel quantization, and evaluating reader studies comparing compact AI assistance with expert neuroradiologist workflows.

